A Key Encapsulation Mechanism (KEM) is an asymmetric encryption technique 
designed to generate both a randomly chosen symmetric key $k_s$ and its 
encrypted version $c$, denoted as {\em encapsulated key}. 
Considering an insecure communication channel, the encapsulation $c$ serves as 
a secure delivery of the secret key, while $k_s$ is meant to be used for the 
subsequent symmetric encryption of data exchanged between the endpoints of the 
communication.
Unlike public key encryption primitives (PKEs), such as either 
RSA~\cite{DBLP:journals/cacm/RivestSA78} or the 
ElGamal~\cite{DBLP:journals/tit/Elgamal85} encryption algorithms, to prepare a proper plaintext there is no need for a criterion to map bitstrings into 
elements of an algebraic structure,
Furthermore, since the bulk of the data is encrypted with the conveyed 
symmetric key $k_s$, they can be received as a stream, without the need to 
buffer the whole ciphertext before decryption. 
From a security standpoint, the only (weak) requirement on the symmetric-key 
cipher used for subsequent data bulk encryption is that it exhibits a 
stream-based mode of operation (e.g., CTR - CounTeR mode) that is secure in case $k_s$ is employed only once per data exchange session.

Post-quantum cryptographic primitives define PKEs that are only weakly secure, 
due to the presence of probabilistic correctness errors or because they are 
proven to exhibit only the computational property of either {\em indistinguishability of ciphertexts under chosen-plaintext attack} (IND-CPA) or 
{\em one-wayness under chosen-plaintext attack} (OW-CPA).
A PKE benefits from OW-CPA if an adversary equipped with a public key and the corresponding encryption of a random plaintext, is computationally unable to 
derive the entire plaintext.
It is well known that IND-CPA implies OW-CPA, because a PKE benefits from the 
IND-CPA property when an adversary is unable to distinguish the encryptions of 
two messages of his choice; more precisely, given the public key and the 
corresponding encryption of one of two attacker-chosen plaintexts, the 
attacker is able to determine which of the two messages was encrypted only 
with negligible success probability. 
%
To guarantee security in both the Random Oracle Model (ROM) and the Quantum Random 
Oracle Model (QROM), every post-quantum cryptographic scheme makes use of a variant 
of the so-called Fujisaki-Okamoto (FO) 
transformation~\cite{DBLP:conf/crypto/FujisakiO99}, introduced for the first 
time in $1999$.
Nowadays, a newer fine-grained modular toolkit of improved FO 
transformations~\cite{DBLP:conf/tcc/HofheinzHK17} is employed for turning weakly 
secure PKE schemes into strongly secure ones, i.e., schemes benefiting from the 
property of {\em indistinguishability of ciphertexts under adaptive chosen-
ciphertext attack} (IND-CCA), i.e., an adversary is unable to distinguish the 
encryptions of two messages of his choice even if he is allowed access to a 
decryption oracle before and after seeing the (randomly selected) ciphertext of one 
of the two messages he has chosen (honoring the constraint of not querying the 
oracle with the challenge ciphertext itself).
\vspace{-0.2cm}
\subsection{HQC cryptosystem}
The first code-based PKE was introduced by McEliece in $1978$ and is based on the 
assumption that the underlying forward error-correction code, a Goppa code, used 
is indistinguishable from a random one. Although the original McEliece 
cryptosystem remains secure to this day as demonstrated by the 
highly optimized Classic-McEliece cryptoscheme~\cite{ClassicMcEliece}, the storage 
requirements for its public key, which is not smaller than $0.25$ MiB for $128$-bit 
security, is deemed too large for many application scenarios. Attempts to reduce 
the key size of the McEliece PKE by using codes more structured than Goppa codes have been all cryptanalyzed.
Therefore, the authors of the IND-CCA {\em Hamming Quasi-Cyclic} (HQC) code-based 
post-quantum KEM proposed a novel approach for the design of their PKE, which 
combines two different types of codes:
{\em (i)} An [$n, k$] code $\mathcal{C}$ defined as the concatenation of 
({\em shortened}) Reed-M\"{u}ller (RM) and ({\em duplicated}) Reed-Solomon (RS) codes, 
for which an efficient decoding algorithm is known, with a fixed, publicly known 
generator matrix $\mathbf{G} \in \mathbb{F}_2^{k \times n}$ ensuring the capability 
to correct up to $\delta$ errors (i.e., bit-flips) in a $n$-bit codeword obtained 
from a message $\mathbf{m} \in \mathbb{F}_2^k$ encoded as a (row) binary vector, as 
$\mathbf{m}\mathbf{G}$. 
{\em (ii)} A random double-circulant [$2n, n$] code with a publicly
known parity-check matrix $\mathbf{H}$$\in$$\mathbb{F}_2^{n \times 2n}$, 
$\mathbf{H}$$=$$[\mathbf{I_n} | \mathbf{rot}(\mathbf{h})]$, where $\mathbf{I_n}$ 
is the identity matrix and $\mathbf{rot}(\mathbf{h}$) is the circulant 
$n \times n$ matrix induced by the $1 \times n$  binary vector $\mathbf{h}$, 
i.e., the matrix obtained from the juxtaposition of $\mathbf{h}$ as the first 
column and the subsequent columns obtained through a vertical rotation, by one 
cell, of the previous one. 
The coefficients of the $1 \times n$  binary vector $\mathbf{h} \in \mathbb{F}_2^n$ are seen as the coefficients of a randomly chosen binary polynomial in 
$\mathcal{R}=\mathbb{F}_2[x]/(x^n-1)$, where $n$ is a prime number such that 
$\mathtt{ord}_2(n) = n-1$, i.e., $2^{n-1} = 1 \bmod n$.
%
The design rationale is that the double-circulant code is going to be used to 
add some ``noise'' (i.e., bit-flips) to the codewords of the inner RM-RS code, 
which in turn will be partially removed through the use of the secret key. The 
leftover noise plus an additional uncorrelated noise term will be decoded by the 
code $\mathcal{C}$, following the same generic steps of the ElGamal 
cryptosystem~\cite{DBLP:journals/tit/Elgamal85} and adapting to quasi-cyclic (QC) 
codes the approach first introduced in~\cite{DBLP:conf/focs/Alekhnovich03}. In such 
a way, HQC can exhibit significantly smaller keys than the Classic-McEliece KEM, 
while still providing the same computational security guarantees. Indeed, for a 
security level of $128$-bit, the public key size is around $2$ kiB vs $256$ kiB 
of the Classic-McEliece scheme.
Specifically, the PKE underlying the HQC design is defined by the following three algorithms.

\noindent \textsf{Keygen()}: samples $\mathbf{h} \xleftarrow{\mathdollar} \mathcal{R}$,  
$(\mathbf{x}, \mathbf{y}) \xleftarrow{\mathdollar} \mathcal{R}_\omega \times \mathcal{R}_\omega$, and 
outputs the private and public keypair as: \kpri = $(\mathbf{x}, \mathbf{y})$ 
and \kpub = $(\mathbf{h}, \mathbf{s} = \mathbf{H}[\mathbf{x}|\mathbf{y}]^{T} = \mathbf{x} + \mathbf{h\cdot y})$. $\mathcal{R}_\omega$ denotes the set of binary polynomials (i.e., vectors in $\mathbb{F}_2^n$) with Hamming Weight $\omega$ (i.e., the number of non-zero coefficients), while $\cdot$ is the circular convolution of the two sequences of coefficients of the factors (each of which is considered as either a binary vector or a binary polynomial); e.g., given $\mathbf{a}, \mathbf{b}, \mathbf{c}=\mathbf{a}\cdot \mathbf{b} \in \mathbb{F}_2^n$, it holds that $\mathbf{a}$$=$$(a_0, \ldots, a_i,\ldots, a_{n-1})$,  $\mathbf{b}$$=$$(b_0, \ldots, b_j,\ldots, b_{n-1})$ and $\displaystyle \mathbf{c} = \mathbf{a} \cdot \mathbf{b} = \hspace{-2.5em}\sum_{\footnotesize \begin{array}{c} k=0 \\ i+j \equiv k \bmod n\end{array}}^{n-1} \hspace{-2.5em} a_ib_j$.

\noindent \textsf{Encrypt(\kpub, \textbf{m})}: takes as input the public key, \kpub = $(\mathbf{h}, \mathbf{s} = \mathbf{x} + \mathbf{h\cdot y})$, samples an ephemeral $1 \times n$ error vector $\mathbf{e}$ from the subset of $\mathcal{R}$ that includes items with Hamming Weight $\omega_e$, i.e., $\mathbf{e} \xleftarrow{\mathdollar} \mathcal{R}_{\omega_e}$, samples the ephemeral vector $r=[\mathbf{r_1}| \mathbf{r_2}] \in \mathbb{F}_2^{2n}$ in such a way that each pair ($\mathbf{r_1}$, $\mathbf{r_2}$) is picked from the Cartesian product of two instances of the subset of $\mathcal{R}$ that includes items with Hamming Weight  $\omega_r$, i.e., $(\mathbf{r_1}, \mathbf{r_2}) \xleftarrow{\mathdollar} \mathcal{R}_{\omega_r} \times \mathcal{R}_{\omega_r}$, and output a ciphertext $\mathbf{c}$ which, in turn, is composed by the pair $(\mathbf{u}, \mathbf{v})$, where: 
$$\mathbf{u} = \mathbf{r_1} + \mathbf{h}\cdot\mathbf{r_2},\ \text{and}\ \mathbf{v} = \mathbf{mG} + \mathbf{s\cdot r_2} + \mathbf{e}$$

\noindent \textsf{Decrypt(\kpri, $(\mathbf{u}, \mathbf{v})$)}: takes as input the private key, \kpri = $(\mathbf{x}, \mathbf{y})$ and a ciphertext $(\mathbf{u}, \mathbf{v}) \in \mathbb{F}_2^n \times \mathbb{F}_2^n$, and outputs the original plaintext message through computing $\mathbf{m} = \mathcal{C}.\textsf{Decode}(\mathbf{v} - \mathbf{u\cdot y})$.

\smallskip
\noindent The computational security of recovering either the secret-key or the plaintext of an instance of the HQC.PKE scheme is computationally equivalent to solving the syndrome decoding problem on random quasi-cyclic codes with length $n$ and dimension $k$~\cite{DBLP:journals/tit/Aguilar-Melchor18}. 
While random quasi-cyclic codes are not proven to yield an NP-complete and hard 
to approximate syndrome decoding problems, it is largely believed that, in case 
$n$ is chosen as a prime such that $\mathtt{ord}(2)$$=$$n$$-$$1$, no speedup 
can be achieved with respect to the solution of the syndrome decoding problem
on random codes with the same length, when the error generating the syndrome 
is half of the code distance expected for a random code~\cite{DBLP:journals/tit/Aguilar-Melchor18}.

Thus, the best classical and quantum cryptanalytic techniques that are currently 
known to perform a key- or a message-recovery attack against the cryptosystem 
have been considered to design the parameters sets shown in Table~\ref{tab:params}. 
In particular, the NIST call defines the security levels for which parameter
sets should be designed stating that the resulting computational effort to 
break the cryptoscheme should match the one to break both classical and quantum realizations of AES-$128$, AES-$192$, and AES-$256$, respectively.
The sets of parameters of a post-quantum cryptosystem that match the security 
desiderata are labeled by NIST as Category $1$, Category $3$ and Category $5$, 
respectively, while the authors of HQC.KEM named the corresponding instances of 
their proposal as \texttt{hqc-128}, \texttt{hqc-192} and \texttt{hqc-256}.
% 
Finally, the most up-to-date specification of HQC reports a public code that is 
defined as the concatenation of Reed-Solomon (outer) and Reed-M\"{u}ller (inner) 
codes, with codeword length equal to $n_1$ and $n_2$, respectively, while the overall 
codeword length $n$ is taken as the smallest prime number greater than $n_1n_2$ 
such that $\mathtt{ord}_n(2)$$=$$n$$-$$1$. Thus, the public generator matrix 
$\mathbf{G}$ has dimension $k\times n$, where $k$ coincides with the bitsize of 
the plaintext message $\mathbf{m}$. Furthermore, the random code with 
parity-check matrix $\mathbf{H}$ is instantiated by randomly sampling the 
vector $\mathbf{h}$ from $\mathds{F}_2^n$. 
\begin{table}[!t]
    \centering
    \caption{HQC Parameter sets. $n$: public code length, $n_1$: Reed–Solomon code length, $n_2$: Reed–M\"{u}ller code length. $\omega$: weight of the $n$-bit binary vectors $\mathbf{x}, \mathbf{y}$, composing the secret key. 
    $\omega_r=\omega_e$: weight of the $n$-bit binary vectors used in the 
    encryption. $[n, k, d]$ for Reed–Solomon and Reed–M\"{u}ller codes are the 
    length, dimension, and minimum distance of the codes, respectively.
    The multiplicity of duplicated Reed–M\"{u}ller code is $3$, $5$, $5$ for \textsc{hqc-128}, \textsc{hqc-192}, \textsc{hqc-256}.
    The probability of a decapsulation error (i.e., of a correctness error of the system) is specified as the decoding failure rate (DFR) of the public code}
    {\small
     \begin{tabular}{cc|cccc|cc|c}
        \toprule
        \textbf{NIST} & \multirow{2}{*}{\textbf{Instance}} & \textbf{Plaintex} & \multirow{1}{*}{\bf 
        Reed-} & \multirow{1}{*}{\bf Reed-} & \multirow{2}{*}{$\mathbf{n}$} & \multirow{2}{*}{$\mathbf{\omega}$} & \multirow{2}{*}{$\mathbf{\omega_r = \omega_e}$} & \multirow{2}{*}{\textbf{DFR}}\\
        \textbf{Category} &  & \textbf{size (bit)} & {\bf Solomon} & {\bf M\"{u}ller}&  &  &  & \\
        \midrule
        $1$& \texttt{hqc-$128$} & $128$ &$[46, 16, 15]$ &  $[384, 8, 192]$ & $17,669$ & $66 $ & $ 75$ & $2^{-128}$\\
        $3$& \texttt{hqc-$192$} & $192$ &$[56, 24, 16]$ &  $[640, 8, 320]$ & $35,851$ & $100$ & $114$ & $2^{-192}$\\
        $5$& \texttt{hqc-$256$} & $256$ &$[90, 32, 29]$ &  $[640, 8, 320]$ & $57,637$ & $131$ & $149$ & $2^{-256}$\\
        \bottomrule
    \end{tabular}
    \label{tab:params}
    }\vspace{-0.4cm}
\end{table}
